\section{Experimental Setup}
\label{sec:setup}

\begin{figure}[t]
  \centering
  % Fig. 1 - experimental pipeline (TikZ). Shared by the report and the slides.
\begin{adjustbox}{max width=\linewidth, max totalheight=0.5\textheight}
\begin{tikzpicture}[
  node distance=5mm and 5mm,
  box/.style={draw=black!55, rounded corners=2pt, fill=#1, align=center, font=\footnotesize,
              minimum height=11mm, text width=22mm, inner sep=2pt},
  box/.default=black!4,
  arr/.style={-{Stealth[length=2mm]}, thick, black!60},
  lab/.style={font=\scriptsize\itshape, text=black!60}]
  % row 1: data
  \node[box] (dl) {Download\\NSMC, Shopping,\\Steam};
  \node[box, right=of dl] (nm) {Normalization\\NFC, cap repeats,\\no NFKC};
  \node[box, right=of nm] (cl) {Cleaning\\empty, conflicts,\\dups};
  \node[box, right=of cl] (sp) {Fixed splits\\8k / 2k / 2k\\(seed 2026)};
  \node[box, right=of sp] (sub) {Nested subsets\\500 $\subset$ \dots $\subset$ 8k\\$\times$ 5 seeds};
  % row 2: modelling
  \node[box=green!8, below=9mm of sub] (tok) {Tokenization\\Kiwi\\morphemes (main)};
  \node[box=blue!8, left=of tok] (tune) {Tuning at $n$=2k\\10 configs $\times$ 9\\(val macro-F1)};
  \node[box=blue!8, left=of tune] (frz) {Frozen configs\\per dataset \&\\method};
  \node[box=blue!8, left=of frz] (fit) {Final fits: 675\\3$\times$9$\times$5$\times$5\\same review IDs};
  \node[box=orange!12, left=of fit] (ev) {Test evaluation\\+ cost (CPU \&\\accelerator)};
  % row 3: ablation branch
  \node[box=purple!8, below=7mm of frz, text width=46mm] (abl) {Tokenizer ablation: best TF-IDF + best neural\\
    $\times$ \{eojeol, Kiwi, char\} $\times$ \{500, 8k\} $\times$ 5 seeds (+120 fits)};
  \draw[arr] (dl) -- (nm);
  \draw[arr] (nm) -- (cl);
  \draw[arr] (cl) -- (sp);
  \draw[arr] (sp) -- (sub);
  \draw[arr] (sub) -- (tok);
  \draw[arr] (tok) -- (tune);
  \draw[arr] (tune) -- (frz);
  \draw[arr] (frz) -- (fit);
  \draw[arr] (fit) -- (ev);
  \draw[arr, dashed] (frz) -- (abl);
  \draw[arr, dashed] (abl.west) -| (ev.south);
\end{tikzpicture}
\end{adjustbox}

  \caption{Experimental pipeline. Validation data drive tuning, early stopping, the choice
  of the best method per family and of the ablation methods; test data are used once, for
  the final evaluation of frozen pipelines. The dashed branch is the tokenizer ablation.}
  \label{fig:pipeline}
\end{figure}

\paragraph{Protocol.}
\Cref{fig:pipeline} summarises the protocol. Every final condition
(3 datasets $\times$ 13 methods $\times$ 10--14 sizes per dataset) is repeated with \nSeeds{} seeds
(\seedList), giving \nFinalFits{} final fits plus \nTuningFits{} tuning fits.
The seed controls the nested training subset and all model randomness, so the
TF-IDF and neural runs with the same seed see exactly the same reviews.

\paragraph{Metrics.}
Macro-F1 is the primary metric; accuracy, macro-precision and macro-recall are
secondary (\cref{tab:a1acc,tab:a1p,tab:a1r}). We report mean\,$\pm$\,sample standard deviation over seeds; the seeds
measure variability but do not constitute a significance test. Confusion matrices at
$n=8{,}000$ are summed over the \nSeeds{} runs and row-normalised. The marginal value of
labels is the per-seed macro-F1 gain between adjacent sizes.

\paragraph{Crossover threshold.}
At every size we select the best TF-IDF and the best neural method by mean
\emph{validation} macro-F1, compute the test macro-F1 gap (neural minus TF-IDF)
per seed---paired, because both saw the same subset---and average it. The
threshold is the smallest size whose mean gap is at least \crossoverMargin{} point;
otherwise we report ``no observed crossover up to the full training set''. As a sensitivity check we
repeat the procedure with winners selected by test macro-F1 (\cref{tab:a4}). At the
threshold (or at the cost-pass size 8{,}000 if there is none) we check whether the neural winner is
Pareto-efficient against the TF-IDF winner on test macro-F1, CPU end-to-end latency,
fit time and peak RAM.

\paragraph{Cost measurement.}
All runs use one machine (\cpuModel{}; accelerator: \gpuModel; \hardwareNote). Every run
is a separate process with a fixed thread count: \nThreads{} CPU threads for TF-IDF fits
and all CPU inference, and \neuralThreads{} for neural training. We run in two modes. In
\emph{quality} runs, which supply all predictive metrics, jobs run concurrently: up to
\kGpu{} neural jobs share the GPU, and \cpuJobs{} TF-IDF jobs run alongside them. Concurrency
changes neither a job's thread count nor its device. In \emph{cost} runs, which supply
every reported time and memory figure, jobs run one at a time on an otherwise idle machine.
These cover all methods at $n=8{,}000$, the crossover sizes and all tuning fits. We check
that the two modes produce the same test predictions (\cref{sec:results}). Neural
models train on \trainDevice{} in \precisionUsed{}; TF-IDF models always run on the CPU. We record end-to-end fit
time (tokenisation + vectoriser/vocabulary + training) and classifier-only fit time;
batch-size-1 inference latency on 200 fixed test reviews after 20 warm-up calls
(median), both classifier-only (pre-vectorised input) and end-to-end (raw string to
label), on the CPU for every method (the primary, comparable measurement) and
additionally on the accelerator (\accelDevice) for neural methods, with device
synchronisation around every call; peak host RAM of the process; peak accelerator
memory; serialised artefact size; and model size
(parameters, non-zero coefficients or tree nodes). Hardware and software versions
are listed in \cref{tab:a6}. As a check on the two-mode design, every cost run was compared
with its quality run: \detIdentical{} of \detRuns{} runs produced identical test
predictions (maximum disagreement \detMaxDisagree), so concurrency did not affect any
predictive result.

\paragraph{Compute used.}
All experiments ran on a single machine with one \gpuModel{} and a \cpuModel{},
in \precisionUsed. Busy wall-clock time per stage, from the first job start to the last job
end, was: pilot \wallPilot{}\,h, tuning \wallTuning{}\,h (exclusive), main sweep
\wallSweep{}\,h (parallel, mean GPU utilisation \gpuUtilSweep), tokenizer ablation
\wallAblation{}\,h (parallel, \gpuUtilAblation) and cost pass \wallCostPass{}\,h
(exclusive). The total is \wallTotal{}\,h, i.e.\ \wallTotal{} GPU-hours on one GPU; the
neural jobs alone add up to \gpuJobHours{}\,h of job time, because up to \kGpu{} of them
shared the GPU. The exclusive stages leave the GPU mostly idle (\gpuUtilTuning{} in tuning,
\gpuUtilCostPass{} in the cost pass), as they must for undisturbed timings.

\paragraph{Deviations from the proposal.}
(i) \nSeeds{} seeds instead of 3, i.e.\ \nFinalFits{} instead of 405 final fits; confusion
matrices are summed over \nSeeds{} runs. Tuning is unchanged (\nTuningFits{} fits).
(ii) A tokenizer ablation (\cref{sec:ablation}, \nAblationFits{} additional fits) that was not
planned. (iii) Preprocessing is specified in more detail than in the proposal (fixed
normalization rules, Kiwi morphemes with POS tags, \cref{sec:preprocessing}).
(iv) All results are reported as mean\,$\pm$\,sample standard deviation over \nSeeds{} seeds.
(v) Cost figures (fit time, latency, memory) are measured in separate exclusive cost runs.
These cover all methods at $n=8{,}000$ and the crossover sizes; they are not measured at every
training size.
(vi) Mamba uses the pure-PyTorch selective scan (\mambaBackend{}), because no
\texttt{mamba-ssm} build exists for the PyTorch version we used. Its fit time and GPU latency are
therefore probably higher than the fused kernels would give.
(vii) Because no crossover was observed, no additional cost runs at a crossover size were
needed; all cost figures refer to $n=8{,}000$.
(viii) The size grid is extended beyond the proposal's maximum of 8{,}000 reviews with
16{,}000, 32{,}000, 64{,}000, steps of 20{,}000 up to 160{,}000 and the full training set
(\cref{sec:data}), so the learning
curves reach the point where all available labels are used. Cost runs stay at
$n=8{,}000$ (plus any crossover size).
(ix) Three recurrent baselines are added: a plain bidirectional Elman RNN (BiRNN), which
separates the effect of recurrence from that of gating in the BiLSTM and BiGRU, and
unidirectional RNN, LSTM and GRU, which isolate the effect of bidirectionality
(\cref{sec:arch}); the proposal's LSTM and GRU are the bidirectional ones.
